% figures08.tex -- Chapter 8 (Parallel Lines) figures redrawn in TikZ.
%
% Local label clearance: ch08 has several figures where a vertex label sits
% at a point that two or three drawn lines pass through, so the shared
% 3.4pt outer sep is not enough.  Raised here only (brumfiel.sty untouched).
\tikzset{
  vlab/.style = {font=\footnotesize\itshape, inner sep=1.2pt, outer sep=5.2pt},
  slab/.style = {font=\scriptsize, inner sep=1.2pt, outer sep=5.2pt},
}
% Geometry read off the photo pages (book pp.133-153 = PDF 147-167) and the
% iPad scans (scan09 pp.7-23 = book pp.133-148; scan10 pp.1-4 = book
% pp.151-153).  Vertex ratios were measured off 400 dpi crops of the source
% figures; constrained points are constructed, never eyeballed: parallels via
% equal ratios, feet via projection, crossings via (intersection of ...).

% ---------------------------------------------------------------- Fig 8-1
% Two rays with the same direction: r1 || r2, both on the same side of the
% line through their vertices (Definition 8-2).
\newcommand{\FIGVIIIONE}{\begin{bkfigure}[0.62\linewidth]
\begin{tikzpicture}[scale=0.9]
  \coordinate (V1) at (0.62,1.15);          % vertex of r1
  \coordinate (V2) at (0.16,0.0);           % vertex of r2
  \draw[aux] ($(V2)!-0.30!(V1)$) -- ($(V2)!1.55!(V1)$);   % line of vertices
  \draw[fig,-latex] (V1) -- ++(4.35,0);
  \draw[fig,-latex] (V2) -- ++(4.35,0);
  \node[slab,above] at ($(V1)+(2.5,0)$) {$r_1$};
  \node[slab,above] at ($(V2)+(2.5,0)$) {$r_2$};
\end{tikzpicture}
\figcap{8--1}
\end{bkfigure}}

% ----------------------------------------------------------- Fig 8-2, 8-3
% 8-2: first proof of Thm 8-1.  n = AB; C on n beyond A; D on m and E on l
% chosen so that angle DAC is congruent to angle EBA, hence m || l, and the
% lines cannot meet at P.  Both congruent angles are arced, as in the book.
% 8-3: the lemma's figure -- the same three lines with m NOT parallel to l,
% and the two corresponding angles numbered 1 and 2.
\newcommand{\FIGVIIITWOTHREE}{\begin{bkfigure}[\linewidth]
\begin{minipage}{0.48\linewidth}\centering
\begin{tikzpicture}[scale=0.88]
  \coordinate (B) at (1.1470,0.0000);      \coordinate (A) at (1.9220,1.4756);
  \coordinate (C) at (2.1993,2.0035);      % on n, beyond A
  \coordinate (D) at (3.0194,1.4756);      % on m
  \coordinate (E) at (2.5110,0.0000);      % on l
  \coordinate (P) at (5.3320,0.7626);
  \draw[fig] (0.4650,0) -- (3.8440,0);                        % l (solid part)
  \draw[fig] (0.9378,-0.3984) -- (2.3295,2.2514);             % n
  \draw[aux] (0.4340,1.4756) -- (3.3790,1.4756);              % m (dashed)
  \draw[aux] (3.3790,1.4756) .. controls +(0.62,-0.06) and +(-0.50,0.38) .. (P);
  \draw[aux] (3.8440,0)      .. controls +(0.72,0.06) and +(-0.66,-0.34) .. (P);
  \draw[fig] ($(P)+(0.16,0.20)$) -- (P) -- ($(P)+(0.16,-0.20)$);
  % the two congruent angles of the construction
  \draw[tick] ($(A)+(0:0.30)$)  arc (0:62.29:0.30);
  \draw[tick] ($(B)+(0:0.36)$)  arc (0:62.29:0.36);
  \foreach \p in {C,D,E}{\fill (\p) circle (1.0pt);}
  \node[vlab] at (1.5100,1.8700) {$A$};
  \node[vlab,right] at (C) {$C$};
  \node[vlab,above] at (D) {$D$};
  \node[vlab,above] at (E) {$E$};   % book sets E above its dot, as with D
  \node[vlab] at (1.5200,-0.3300) {$B$};
  \node[vlab,right] at (P) {$P$};
  \node[vlab,below right] at (D) {$m$};
  \node[vlab,above] at (1.75,0) {$l$};
  \node[vlab] at (1.1200,0.8500) {$n$};
\end{tikzpicture}
\figcap{8--2}
\end{minipage}\hfill
\begin{minipage}{0.48\linewidth}\centering
\begin{tikzpicture}[scale=0.88]
  \coordinate (A) at (1.5190,1.4756);   \coordinate (B) at (0.7750,0.0000);
  \draw[fig] (-0.0250,-0.0594) -- (3.1950,0.1796);            % l through B
  \draw[fig] (0.7990,1.6310)  -- (3.5690,1.0332);             % m through A
  \draw[fig] (0.5769,-0.3929) -- (1.7531,1.9399);             % n through A,B
  \draw[fig,-latex] (2.387,0.620) -- (2.945,0.620);
  \node[vlab,right] at (2.945,0.620) {$P$};
  \draw[tick] ($(A)+(-12.18:0.34)$) arc (-12.18:63.24:0.34);
  \draw[tick] ($(B)+(4.24:0.38)$)   arc (4.24:63.24:0.38);
  \node[slab] at (2.0604,1.7342) {$1$};
  \node[slab] at (1.3238,0.3666) {$2$};
  \node[vlab] at (0.9050,1.1830) {$A$};
  \node[vlab] at (0.1400,-0.4000) {$B$};
  \node[vlab,right] at (3.5690,1.0332) {$m$};
  \node[vlab,above] at (2.1000,0.0983) {$l$};
  \node[vlab] at (0.5000,0.7800) {$n$};
\end{tikzpicture}
\figcap{8--3}
\end{minipage}
\end{bkfigure}}

% ---------------------------------------------------------------- Fig 8-4
% Third proof: n cuts m at A and l at B, with the indicated angles at A and B
% congruent (both arced).  Q is on m with AQ congruent to BP.
% The book draws m and l PARALLEL (both horizontal) and bends each with a
% dashed curve to the hypothetical meeting point P -- the same convention as
% Fig 8-2.  That is what makes the two arced angles at A and B congruent, as
% the third proof requires; m must not be tilted off l.
\newcommand{\FIGVIIIFOUR}{\begin{bkfigure}[0.66\linewidth]
\begin{tikzpicture}[scale=0.86]
  \coordinate (B) at (2.0460,0.0000);  \coordinate (A) at (2.8830,1.6244);
  \coordinate (P) at (4.5880,0.9424);
  % Q on m, with m || l through A and AQ = BP = 2.7111 exactly
  % (both re-checked in tools/constraints/ch08.py)
  \coordinate (Q) at (0.1719,1.6244);
  \draw[fig] (0.0000,0) -- (B);                          % l, drawn up to B
  \draw[fig] (1.9269,-0.2311) -- (3.0342,1.9177);        % n
  \draw[aux] (Q) -- (A);                                 % m
  \draw[fig] (Q) -- (B);
  \draw[aux] (A) .. controls +(1.00,-0.07) and +(-0.52,0.40) .. (P);
  \draw[aux] (B) .. controls +(1.08,0.10) and +(-0.66,-0.42) .. (P);
  \draw[fig] ($(P)+(0.17,0.20)$) -- (P) -- ($(P)+(0.17,-0.20)$);
  % the two congruent angles at A and B (both 62.74 deg, since m || l)
  \draw[tick] ($(A)+(0:0.32)$) arc (0:62.74:0.32);
  \draw[tick] ($(B)+(0:0.36)$) arc (0:62.74:0.36);
  \node[vlab,above left] at (A) {$A$};
  \node[vlab,below right] at (B) {$B$};
  \node[vlab,above] at (Q) {$Q$};
  \node[vlab,right] at (P) {$P$};
  \node[vlab,left] at (Q) {$m$};
  \node[vlab,left] at (0.0000,0) {$l$};
  \node[vlab,above right] at (3.0342,1.9177) {$n$};
\end{tikzpicture}
\figcap{8--4}
\end{bkfigure}}

% ----------------------------------------------------------- Fig 8-5, 8-6
% 8-5: l || m and m || n -- all three drawn parallel -- with l and n bent
% (dashed) to meet at P, the contradiction.
% 8-6: l || m and transversal n through P on l.
\newcommand{\FIGVIIIFIVESIX}{\begin{bkfigure}[\linewidth]
\begin{minipage}{0.48\linewidth}\centering
\begin{tikzpicture}[scale=0.80]
  \draw[fig] (0.30,1.90) -- (2.55,1.90);                  % l
  \draw[fig] (0.30,1.12) -- (2.55,1.12);                  % n
  \draw[fig] (0.30,0.34) -- (2.55,0.34);                  % m
  \coordinate (P) at (3.60,1.50);
  \draw[aux] (2.55,1.90) -- (P);
  \draw[aux] (2.55,1.12) -- (P);
  \draw[fig] ($(P)+(0.17,0.20)$) -- (P) -- ($(P)+(0.17,-0.20)$);
  \node[vlab,left] at (0.30,1.90) {$l$};
  \node[vlab,left] at (0.30,1.12) {$n$};
  \node[vlab,left] at (0.30,0.34) {$m$};
  \node[vlab,right] at (P) {$P$};
\end{tikzpicture}
\figcap{8--5}
\end{minipage}\hfill
\begin{minipage}{0.48\linewidth}\centering
\begin{tikzpicture}[scale=0.80]
  \draw[fig] (0.20,1.70) -- (3.70,1.70);                  % l
  \draw[fig] (0.20,0.42) -- (3.70,0.42);                  % m
  \coordinate (P) at (1.72,1.70);
  \draw[fig] ($(P)!-0.28!(2.86,0.10)$) -- (2.86,0.10);    % n
  \fill (P) circle (1.0pt);
  \node[vlab,above right] at (P) {$P$};
  \node[vlab,right] at (3.70,1.70) {$l$};
  \node[vlab,right] at (3.70,0.42) {$m$};
  \node[vlab,below] at (2.86,0.10) {$n$};
\end{tikzpicture}
\figcap{8--6}
\end{minipage}
\end{bkfigure}}

% ----------------------------------------------------------- Fig 8-7, 8-8
% 8-7: angle ABC congruent to angle BCD, hence AB || CD.
% 8-8: AB = CD, angles at B and C right, hence AD || BC.
\newcommand{\FIGVIIISEVENEIGHT}{\begin{bkfigure}[\linewidth]
\begin{minipage}{0.50\linewidth}\centering
\begin{tikzpicture}[scale=0.80]
  \coordinate (A) at (0,0);    \coordinate (B) at (1.72,0);
  \coordinate (C) at (2.52,0.86);
  \coordinate (D) at ($(C)+(1.72,0)$);
  \draw[fig] (A) -- (B) -- (C) -- (D);
  \node[vlab,left] at (A) {$A$};
  \node[vlab,below right] at (B) {$B$};
  \node[vlab,above left] at (C) {$C$};
  \node[vlab,right] at (D) {$D$};
\end{tikzpicture}
\figcap{8--7}
\end{minipage}\hfill
\begin{minipage}{0.46\linewidth}\centering
\begin{tikzpicture}[scale=0.80]
  \coordinate (B) at (0,0);    \coordinate (C) at (3.05,0);
  \coordinate (A) at (0,1.72); \coordinate (D) at (3.05,1.72);
  \draw[fig] (A) -- (B) -- (C) -- (D) -- cycle;
  % right-angle boxes at B and C
  \draw[fig] ($(B)+(0,0.20)$) -- ($(B)+(0.20,0.20)$) -- ($(B)+(0.20,0)$);
  \draw[fig] ($(C)+(0,0.20)$) -- ($(C)+(-0.20,0.20)$) -- ($(C)+(-0.20,0)$);
  % double ticks on AB and DC
  \foreach \dy in {-0.07,0.07}{
    \draw[tick] ($(A)!0.5!(B)+(-0.13,\dy)$) -- ($(A)!0.5!(B)+(0.13,\dy)$);
    \draw[tick] ($(D)!0.5!(C)+(-0.13,\dy)$) -- ($(D)!0.5!(C)+(0.13,\dy)$);}
  \node[vlab,above left] at (A) {$A$};
  \node[vlab,below left] at (B) {$B$};
  \node[vlab,below right] at (C) {$C$};
  \node[vlab,above right] at (D) {$D$};
\end{tikzpicture}
\figcap{8--8}
\end{minipage}
\end{bkfigure}}

% --------------------------------------------------------- Fig 8-9, 8-10
% 8-9: transversal l cutting l1 and l2; the eight numbered angles.  As in the
% book, l leans the other way from l1/l2 -- its upper crossing is to the LEFT
% of its lower one.  Each number sits on the bisector of its own angle.
% 8-10: proof figure for Thm 8-4, with G chosen so that AG || l2 (drawn just
% off l1 on purpose -- the proof concludes that G must in fact lie on l1).
\newcommand{\FIGVIIININETEN}{\begin{bkfigure}[\linewidth]
\begin{minipage}{0.50\linewidth}\centering
\begin{tikzpicture}[scale=1.30]
  \coordinate (L1a) at (-0.0620,1.1036);  \coordinate (L1b) at (3.0380,1.0745);
  \coordinate (L2a) at ( 0.0620,0.1240);  \coordinate (L2b) at (3.0380,0.3887);
  \coordinate (La)  at ( 1.3640,1.5047);  \coordinate (Lb)  at (2.3355,0.0000);
  \draw[fig] (L1a) -- (L1b);                              % l1
  \draw[fig] (L2a) -- (L2b);                              % l2
  \draw[fig] (La) -- (Lb);                                % l
  \node[slab] at (1.8070,1.4030) {$1$};
  \node[slab] at (1.3180,1.2614) {$2$};
  \node[slab] at (1.4596,0.7724) {$3$};
  \node[slab] at (1.9486,0.9140) {$4$};
  \node[slab] at (2.2943,0.6320) {$5$};
  \node[slab] at (1.8128,0.4665) {$6$};
  \node[slab] at (1.9783,-0.0150) {$7$};
  \node[slab] at (2.4598,0.1505) {$8$};
  \node[vlab,right] at (L1b) {$l_1$};
  \node[vlab,right] at (L2b) {$l_2$};
  \node[vlab,above] at (La) {$l$};
\end{tikzpicture}
\figcap{8--9}
\end{minipage}\hfill
\begin{minipage}{0.46\linewidth}\centering
\begin{tikzpicture}[scale=0.92]
  \coordinate (A) at (1.62,1.42);  \coordinate (B) at (1.34,0.30);
  \draw[fig] (-0.20,1.42) -- (3.60,1.42);                 % l1
  \draw[fig] (-0.20,0.30) -- (3.60,0.30);                 % l2
  \coordinate (E) at ($(B)!1.68!(A)$);
  \draw[fig] ($(B)!-0.42!(A)$) -- (E);                    % l
  \coordinate (C) at (0.42,1.42);
  \coordinate (F) at (3.00,1.42);
  \coordinate (D) at (2.85,0.30);
  % G on the ray from A that the proof compares with l1
  \coordinate (G) at (2.60,1.86);
  \draw[aux] (A) -- (G);
  \foreach \p in {C,F,D,G,E}{\fill (\p) circle (1.0pt);}  % book dots E too
  \node[vlab,above] at (C) {$C$};
  \node[vlab] at (1.374,1.735) {$A$};
  \node[vlab,above] at (F) {$F$};
  \node[vlab] at (1.586,-0.015) {$B$};
  \node[vlab,above] at (D) {$D$};
  \node[vlab,above left] at (E) {$E$};
  \node[vlab,above left] at (G) {$G$};
  \node[vlab,right] at (3.60,1.42) {$l_1$};
  \node[vlab,right] at (3.60,0.30) {$l_2$};
  \node[vlab,below left] at ($(B)!-0.42!(A)$) {$l$};
\end{tikzpicture}
\figcap{8--10}
\end{minipage}
\end{bkfigure}}

% --------------------------------------------------------------- Fig 8-11
% Bowtie: AO = CO and BO = OD, so AB || CD and BC || AD.
\newcommand{\FIGVIIIELEVEN}{\begin{bkfigure}[0.56\linewidth]
\begin{tikzpicture}[scale=0.86]
  \coordinate (O) at (2.05,0.92);
  \coordinate (A) at (1.15,0);      \coordinate (C) at ($(A)!2!(O)$);
  \coordinate (B) at (0,0.72);      \coordinate (D) at ($(B)!2!(O)$);
  \draw[fig] (A) -- (C);
  \draw[fig] (B) -- (D);
  \draw[fig] (B) -- (A);
  \draw[fig] (C) -- (D);
  \node[vlab,below] at (A) {$A$};
  \node[vlab,left]  at (B) {$B$};
  \node[vlab,above] at (C) {$C$};
  \node[vlab,right] at (D) {$D$};
  \node[vlab] at (1.86,1.36) {$O$};
\end{tikzpicture}
\figcap{8--11}
\end{bkfigure}}

% --------------------------------------------------------------- Fig 8-12
% Adjacent angles AOB and BOC: OB interior to angle AOC.
\newcommand{\FIGVIIITWELVE}{\begin{bkfigure}[0.48\linewidth]
\begin{tikzpicture}[scale=0.90]
  \coordinate (O) at (0,0);
  \coordinate (A) at (76:2.15);
  \coordinate (B) at (48:1.80);
  \coordinate (C) at (-9:2.45);
  \draw[fig] (O) -- (A);
  \draw[fig] (O) -- (B);
  \draw[fig] (O) -- (C);
  \node[vlab,above] at (A) {$A$};
  \node[vlab,above right] at (B) {$B$};
  \node[vlab,right] at (C) {$C$};
  \node[vlab,left] at (O) {$O$};
\end{tikzpicture}
\figcap{8--12}
\end{bkfigure}}

% --------------------------------------------------------------- Fig 8-13
% Thm 8-5: l through A parallel to BC; D and E on l, one on each side of AB.
\newcommand{\FIGVIIITHIRTEEN}{\begin{bkfigure}[0.58\linewidth]
\begin{tikzpicture}[scale=0.88]
  \coordinate (B) at (0,0);  \coordinate (C) at (3.35,0);
  \coordinate (A) at (1.675,1.90);
  \draw[fig] (B) -- (C) -- (A) -- cycle;
  % l is the horizontal through A: same direction as BC by construction
  \coordinate (D) at ($(A)+(-1.15,0)$);
  \coordinate (E) at ($(A)+(1.15,0)$);
  \draw[fig] ($(A)+(-1.72,0)$) -- ($(A)+(1.72,0)$);
  \fill (D) circle (1.0pt); \fill (E) circle (1.0pt);
  \node[vlab,above] at (D) {$D$};
  \node[vlab,above] at (A) {$A$};
  \node[vlab,above] at (E) {$E$};
  \node[vlab,left]  at (B) {$B$};
  \node[vlab,right] at (C) {$C$};
  \node[vlab,right] at ($(A)+(1.72,0)$) {$l$};
\end{tikzpicture}
\figcap{8--13}
\end{bkfigure}}

% --------------------------------------------------------------- Fig 8-14
% Two triangles with AB congruent to A'B', angle A congruent to angle A',
% angle C congruent to angle C'.  The book sets the primed triangle up and to
% the right, overlapping the first in x, and mirrored in sense.
\newcommand{\FIGVIIIFOURTEEN}{\begin{bkfigure}[0.66\linewidth]
\begin{tikzpicture}[scale=0.95]
  \coordinate (A) at (0,1.35);  \coordinate (B) at (0.39,0);
  \coordinate (C) at (2.34,0.78);
  \draw[fig] (A) -- (B) -- (C) -- cycle;
  \node[vlab,left] at (A) {$A$};
  \node[vlab,below] at (B) {$B$};
  \node[vlab,below right] at (C) {$C$};
  % A'B'C' solved from the three side lengths of ABC, so it is congruent to
  % ABC exactly (opposite sense, as the book draws it); checked in ch08.py.
  \coordinate (Cp) at (1.2000,2.1000);
  \coordinate (Ap) at (3.4950,2.8302);
  \coordinate (Bp) at (3.1992,1.4565);
  \draw[fig] (Ap) -- (Bp) -- (Cp) -- cycle;
  \node[vlab,above right] at (Ap) {$A'$};
  \node[vlab,below right] at (Bp) {$B'$};
  \node[vlab,left] at (Cp) {$C'$};
\end{tikzpicture}
\figcap{8--14}
\end{bkfigure}}

% -------------------------------------------------------- Fig 8-15, 8-16
% 8-15: angle A right, AD perpendicular to BC (D is the foot).
% 8-16: D the midpoint of BC, with BD congruent to AD.
\newcommand{\FIGVIIIFIFTEENSIXTEEN}{\begin{bkfigure}[\linewidth]
\begin{minipage}{0.48\linewidth}\centering
\begin{tikzpicture}[scale=0.80]
  % right angle at A: put A on the circle with diameter BC
  \coordinate (B) at (0,0);  \coordinate (C) at (3.45,0);
  \coordinate (M) at ($(B)!0.5!(C)$);
  \coordinate (A) at ($(M)+(65:1.725)$);
  \coordinate (D) at ($(B)!(A)!(C)$);       % foot of the altitude
  \draw[fig] (B) -- (C) -- (A) -- cycle;
  \draw[fig] (A) -- (D);
  \node[vlab,above] at (A) {$A$};
  \node[vlab,left]  at (B) {$B$};
  \node[vlab,right] at (C) {$C$};
  \node[vlab,below] at (D) {$D$};
\end{tikzpicture}
\figcap{8--15}
\end{minipage}\hfill
\begin{minipage}{0.48\linewidth}\centering
\begin{tikzpicture}[scale=0.80]
  \coordinate (B) at (0,0);  \coordinate (C) at (3.45,0);
  \coordinate (D) at ($(B)!0.5!(C)$);       % midpoint of BC
  % A on the circle centred D through B, so that BD = AD
  \coordinate (A) at ($(D)+(62:1.725)$);
  \draw[fig] (B) -- (C) -- (A) -- cycle;
  \draw[fig] (A) -- (D);
  \node[vlab,above] at (A) {$A$};
  \node[vlab,left]  at (B) {$B$};
  \node[vlab,right] at (C) {$C$};
  \node[vlab,below] at (D) {$D$};
\end{tikzpicture}
\figcap{8--16}
\end{minipage}
\end{bkfigure}}

% -------------------------------------------------------- Fig 8-17, 8-18
% 8-17: a polygonal path P1 ... Pn.  8-18: a six-sided path that crosses itself.
\newcommand{\FIGVIIISEVENTEENEIGHTEEN}{\begin{bkfigure}[\linewidth]
\begin{minipage}{0.48\linewidth}\centering
\begin{tikzpicture}[scale=0.74]
  \coordinate (P1) at (0,0.30);
  \coordinate (P2) at (1.15,0.92);
  \coordinate (P3) at (0.62,2.62);
  \coordinate (Q1) at (1.72,1.72);
  \coordinate (Q2) at (2.35,2.42);
  \coordinate (Q3) at (2.95,1.05);
  \coordinate (Pn) at (2.28,0);
  \draw[fig] (P1) -- (P2) -- (P3) -- (Q1) -- (Q2) -- (Q3) -- (Pn);
  \node[vlab,left]  at (P1) {$P_1$};
  \node[vlab,below right] at (P2) {$P_2$};
  \node[vlab,above left] at (P3) {$P_3$};
  \node[vlab,below right] at (Pn) {$P_n$};
\end{tikzpicture}
\figcap{8--17}
\end{minipage}\hfill
\begin{minipage}{0.48\linewidth}\centering
\begin{tikzpicture}[scale=0.74]
  \coordinate (P1) at (0,0.62);
  \coordinate (P2) at (3.35,1.42);
  \coordinate (P3) at (1.05,2.62);
  \coordinate (P4) at (0.62,0);
  \coordinate (P5) at (1.95,2.62);
  \coordinate (P6) at (2.05,0);
  \coordinate (P7) at (0.30,1.95);
  \draw[fig] (P1) -- (P2) -- (P3) -- (P4) -- (P5) -- (P6) -- (P7);
  \node[vlab,left]  at (P1) {$P_1$};
  \node[vlab,right] at (P2) {$P_2$};
  \node[vlab,above] at (P3) {$P_3$};
  \node[vlab,below] at (P4) {$P_4$};
  \node[vlab,above] at (P5) {$P_5$};
  \node[vlab,below] at (P6) {$P_6$};
  \node[vlab,left]  at (P7) {$P_7$};
\end{tikzpicture}
\figcap{8--18}
\end{minipage}
\end{bkfigure}}

% --------------------------------------------------------------- Fig 8-19
% Three simple polygons: a quadrilateral, a triangle, and a non-convex
% pentagon (the notch at P2).
\newcommand{\FIGVIIININETEEN}{\begin{bkfigure}[\linewidth]
\begin{minipage}{0.31\linewidth}\centering
\begin{tikzpicture}[scale=0.66]
  % measured off the book crop: both uprights vertical, top edge rising to
  % the right (the earlier version had it falling to the right)
  \coordinate (P1) at (0,0);       \coordinate (P4) at (2.350,0.044);
  \coordinate (P3) at (2.350,1.862); \coordinate (P2) at (0,1.703);
  \draw[fig] (P1) -- (P4) -- (P3) -- (P2) -- cycle;
  \node[vlab,below left] at (P1) {$P_1$};
  \node[vlab,above left] at (P2) {$P_2$};
  \node[vlab,above right] at (P3) {$P_3$};
  \node[vlab,below right] at (P4) {$P_4$};
\end{tikzpicture}
\end{minipage}\hfill
\begin{minipage}{0.31\linewidth}\centering
\begin{tikzpicture}[scale=0.66]
  \coordinate (P1) at (0,0);   \coordinate (P3) at (2.62,0);
  \coordinate (P2) at (1.01,1.82);   % apex at 0.386 of the base, as measured
  \draw[fig] (P1) -- (P3) -- (P2) -- cycle;
  \node[vlab,below left] at (P1) {$P_1$};
  \node[vlab,above] at (P2) {$P_2$};
  \node[vlab,below right] at (P3) {$P_3$};
\end{tikzpicture}
\end{minipage}\hfill
\begin{minipage}{0.31\linewidth}\centering
\begin{tikzpicture}[scale=0.66]
  % measured off the book crop: P1P2 nearly level, notch (reflex vertex) at P2
  \coordinate (P1) at (0,1.275);   \coordinate (P2) at (1.214,1.213);
  \coordinate (P3) at (1.565,2.418); \coordinate (P4) at (2.621,1.803);
  \coordinate (P5) at (1.346,0);
  \draw[fig] (P1) -- (P2) -- (P3) -- (P4) -- (P5) -- cycle;
  \node[vlab,left] at (P1) {$P_1$};
  % book sets P2 below-right of the notch; placed explicitly so the glyph
  % clears side P4P5, which runs close by on that side
  \node[vlab] at (1.250,0.700) {$P_2$};
  \node[vlab,above] at (P3) {$P_3$};
  \node[vlab,right] at (P4) {$P_4$};
  \node[vlab,below] at (P5) {$P_5$};
\end{tikzpicture}
\end{minipage}
\figcap{8--19.\ \ Simple polygons.}
\end{bkfigure}}

% -------------------------------------------------------- Fig 8-20, 8-21
% 8-20: naming parts of a polygon -- a quadrilateral with two consecutive
% angles arced, and (as the book draws it) a PENTAGON carrying two of its
% diagonals, dashed, with a pair of adjacent sides called out.
% 8-21: convex and non-convex.
\newcommand{\FIGVIIITWENTYTWENTYONE}{\begin{bkfigure}[\linewidth]
\begin{minipage}{0.60\linewidth}\centering
\begin{tikzpicture}[scale=0.836]
  % left quadrilateral: two consecutive angles marked with arcs
  \coordinate (a1) at (0,0);       \coordinate (a2) at (2.42,0);
  \coordinate (a3) at (2.42,1.42); \coordinate (a4) at (0.55,1.42);
  \draw[fig] (a1) -- (a2) -- (a3) -- (a4) -- cycle;
  \draw[tick] ($(a3)+(180:0.36)$) arc (180:270:0.36);
  \draw[tick] ($(a2)+(180:0.36)$) arc (180:90:0.36);
  \node[slab,align=center] at (1.15,2.70) {Consecutive\\angles};
  \draw[fig,-latex,line width=0.4pt] (1.20,2.00) -- (2.14,1.19);
  \draw[fig,-latex,line width=0.4pt] (1.34,2.00) -- (2.14,0.29);
  % right pentagon: adjacent sides and two diagonals
  \begin{scope}[shift={(3.10,0)}]
    \coordinate (b1) at (0,0);           % bottom-left
    \coordinate (b2) at (0.373,1.779);   % top-left
    \coordinate (b3) at (2.315,1.779);   % top-right
    \coordinate (b4) at (2.929,0.886);   % right point
    \coordinate (b5) at (2.286,0);       % bottom-right
    \draw[fig] (b1) -- (b2) -- (b3) -- (b4) -- (b5) -- cycle;
    \draw[aux] (b1) -- (b3);
    \draw[aux] (b3) -- (b5);
    \node[slab,align=center] at (0.35,2.70) {Adjacent\\sides};
    \node[slab] at (2.45,2.86) {Diagonals};
    \draw[fig,-latex,line width=0.4pt] (0.24,2.00) -- (0.131,0.623);
    \draw[fig,-latex,line width=0.4pt] (0.62,2.00) -- (0.850,1.779);
    \draw[fig,-latex,line width=0.4pt] (2.18,2.40) -- (1.572,1.208);
    \draw[fig,-latex,line width=0.4pt] (2.62,2.40) -- (2.301,0.943);
  \end{scope}
\end{tikzpicture}
\figcap{8--20}
\end{minipage}\hfill
\begin{minipage}{0.36\linewidth}\centering
\begin{tikzpicture}[scale=0.62]
  \coordinate (c1) at (0,0);   \coordinate (c2) at (1.86,0);
  \coordinate (c3) at (2.05,1.42); \coordinate (c4) at (0.42,1.62);
  \draw[fig] (c1) -- (c2) -- (c3) -- (c4) -- cycle;
  \node[slab] at (1.05,2.20) {Convex};
  \begin{scope}[shift={(2.75,0)}]
    \coordinate (d1) at (0,0.55);  \coordinate (d2) at (1.30,1.05);
    \coordinate (d3) at (2.62,0);  \coordinate (d4) at (1.15,0.62);
    \draw[fig] (d1) -- (d2) -- (d3) -- (d4) -- cycle;
    \node[slab] at (1.30,2.20) {Not convex};
  \end{scope}
\end{tikzpicture}
\figcap{8--21}
\end{minipage}
\end{bkfigure}}

% -------------------------------------------------------- Fig 8-22, 8-23
% 8-22: a pentagon cut into three triangles by the two diagonals (dashed)
% from one vertex, with all five exterior angles extended (dashed) and arced.
% 8-23: the same idea for an n-gon -- the book draws a seven-gon split by the
% four diagonals from one vertex, with two exterior angles suggested.
\newcommand{\FIGVIIITWENTYTWOTHREE}{\begin{bkfigure}[\linewidth]
\begin{minipage}{0.48\linewidth}\centering
\begin{tikzpicture}[scale=0.88]
  % Vertex ratios measured off a 400 dpi crop of the book figure with a 50 px
  % grid (book p.147): side v1v2 horizontal, side v2v3 vertical, v4 nearly
  % above v1.  The earlier coordinates were noticeably taller and narrower.
  \coordinate (v1) at (1.3360,0.0000); \coordinate (v2) at (2.9500,0.0230);
  \coordinate (v3) at (2.9090,1.0060); \coordinate (v4) at (1.0290,1.6310);
  \coordinate (v5) at (0.0000,0.5490);
  \draw[fig] (v1) -- (v2) -- (v3) -- (v4) -- (v5) -- cycle;
  \draw[aux] (v1) -- (v3);
  \draw[aux] (v1) -- (v4);
  % exterior-angle extensions: at each vertex, extend the side arrived on
  \draw[aux] (v1) -- (1.7060,-0.1520);
  \draw[aux] (v2) -- (3.3500,0.0287);
  \draw[aux] (v3) -- (2.8923,1.4057);
  \draw[aux] (v4) -- (0.6494,1.7572);
  \draw[aux] (v5) -- (-0.2757,0.2591);
  \draw[tick] ($(v1)+(337.66:0.26)$) arc (337.66:360.82:0.26);
  \draw[tick] ($(v2)+(0.82:0.26)$)   arc (0.82:92.39:0.26);
  \draw[tick] ($(v3)+(92.39:0.26)$)  arc (92.39:161.61:0.26);
  \draw[tick] ($(v4)+(161.61:0.26)$) arc (161.61:226.44:0.26);
  \draw[tick] ($(v5)+(226.44:0.26)$) arc (226.44:337.66:0.26);
  % each number sits on its exterior angle's bisector, at a radius chosen so
  % the glyph clears both bounding rays (the angle at v1 is only 23 deg, so
  % that one has to sit further out)
  \node[slab] at (2.1711,-0.1587) {$1$};
  \node[slab] at (3.2935,0.3863)  {$2$};
  \node[slab] at (2.6081,1.4053)  {$3$};
  \node[slab] at (0.5245,1.5050)  {$4$};
  \node[slab] at (0.1044,0.0600)  {$5$};
\end{tikzpicture}
\figcap{8--22.\ \ Pentagon\\(three triangles).}
\end{minipage}\hfill
\begin{minipage}{0.48\linewidth}\centering
\begin{tikzpicture}[scale=0.88]
  \coordinate (w1) at (0.0000,1.6150);   % the vertex the diagonals come from
  \coordinate (w2) at (0.3610,2.6980);
  \coordinate (w3) at (0.8740,2.7930);
  \coordinate (w4) at (3.0400,1.4630);
  \coordinate (w5) at (2.8690,0.6840);
  \coordinate (w6) at (1.9190,0.0000);
  \coordinate (w7) at (0.7600,0.0380);
  \draw[fig] (w1) -- (w2) -- (w3) -- (w4) -- (w5) -- (w6) -- (w7) -- cycle;
  \foreach \v in {w3,w4,w5,w6}{\draw[fig] (w1) -- (\v);}
  \draw[aux] (w3) -- (1.4836,2.9059);
  \draw[aux] (w4) -- (3.1729,2.0686);
\end{tikzpicture}
\figcap{8--23.\ \ $n$-gon\\($n{-}2$ triangles).}
\end{minipage}
\end{bkfigure}}

% --------------------------------------------------------------- Fig 8-24
% Regular five-pointed star: the points are the vertices of a regular pentagon.
\newcommand{\FIGVIIITWENTYFOUR}{\begin{bkfigure}[0.42\linewidth]
\begin{tikzpicture}[scale=0.90]
  \foreach \i in {0,...,4}{\coordinate (s\i) at ({90+72*\i}:1.72);}
  \draw[fig] (s0) -- (s2) -- (s4) -- (s1) -- (s3) -- cycle;
\end{tikzpicture}
\figcap{8--24}
\end{bkfigure}}

% --------------------------------------------------------------- Fig 8-25
% Parallelogram, rectangle, square, rhombus.
\newcommand{\FIGVIIITWENTYFIVE}{\begin{bkfigure}[\linewidth]
\begin{tikzpicture}[scale=0.70]
  % parallelogram
  \begin{scope}
    \draw[fig] (0,0) -- (1.42,0) -- (1.86,1.72) -- (0.44,1.72) -- cycle;
    \node[slab,below] at (0.92,-0.12) {Parallelogram};
  \end{scope}
  % rectangle
  \begin{scope}[shift={(3.05,0)}]
    \draw[fig] (0,0) -- (2.05,0) -- (2.05,1.62) -- (0,1.62) -- cycle;
    \node[slab,below] at (1.02,-0.12) {Rectangle};
  \end{scope}
  % square
  \begin{scope}[shift={(6.30,0)}]
    \draw[fig] (0,0) -- (1.62,0) -- (1.62,1.62) -- (0,1.62) -- cycle;
    \node[slab,below] at (0.80,-0.12) {Square};
  \end{scope}
  % rhombus
  \begin{scope}[shift={(8.95,0)}]
    \draw[fig] (0,0) -- (1.62,0) -- (2.06,1.56) -- (0.44,1.56) -- cycle;
    \node[slab,below] at (1.02,-0.12) {Rhombus};
  \end{scope}
\end{tikzpicture}
\figcap{8--25}
\end{bkfigure}}

% -------------------------------------------------------- Fig 8-26, 8-27
% 8-26: DE || AC and DF || AB.   8-27: regular hexagon ABCDEF, diagonal FD.
\newcommand{\FIGVIIITWENTYSIXSEVEN}{\begin{bkfigure}[\linewidth]
\begin{minipage}{0.50\linewidth}\centering
\begin{tikzpicture}[scale=0.80]
  \coordinate (B) at (0,0);  \coordinate (C) at (3.45,0);
  \coordinate (A) at (1.72,1.86);
  \coordinate (D) at ($(B)!0.52!(C)$);
  % E on AB and F on AC with DE || AC, DF || AB  (equal ratios)
  \coordinate (E) at ($(B)!0.52!(A)$);
  \coordinate (F) at ($(C)!0.48!(A)$);
  \draw[fig] (B) -- (C) -- (A) -- cycle;
  \draw[aux] (D) -- (E);
  \draw[aux] (D) -- (F);
  \node[vlab,above] at (A) {$A$};
  \node[vlab,left]  at (B) {$B$};
  \node[vlab,right] at (C) {$C$};
  \node[vlab,below] at (D) {$D$};
  \node[vlab,above left]  at (E) {$E$};
  \node[vlab,above right] at (F) {$F$};
\end{tikzpicture}
\figcap{8--26}
\end{minipage}\hfill
\begin{minipage}{0.46\linewidth}\centering
\begin{tikzpicture}[scale=0.80]
  \foreach \i in {0,...,5}{\coordinate (h\i) at ({-60+60*\i}:1.42);}
  \draw[fig] (h0) -- (h1) -- (h2) -- (h3) -- (h4) -- (h5) -- cycle;
  \draw[aux] (h4) -- (h2);
  \node[vlab,below right] at (h0) {$B$};
  \node[vlab,right] at (h1) {$C$};
  \node[vlab,above right] at (h2) {$D$};
  \node[vlab,above left] at (h3) {$E$};
  \node[vlab,left] at (h4) {$F$};
  \node[vlab,below left] at (h5) {$A$};
\end{tikzpicture}
\figcap{8--27}
\end{minipage}
\end{bkfigure}}

% -------------------------------------------------------- Fig 8-28, 8-29
% 8-28: two cevians AQ and AR from A.  The lettered angles are, exactly as
% the book places them:  x = angle B, y = angle BAQ, s = angle RAC,
% z = angle AQC (right of AQ), t = angle ARB (left of AR), r = angle C.
% With that reading x + y + t = r + s + z, which is what Ex. 9 asks for; it
% is verified numerically in tools/constraints/ch08.py.
% 8-29: AB = BC, BE = BD, angle CBD = 40 degrees; angle B obtuse as drawn.
\newcommand{\FIGVIIITWENTYEIGHTNINE}{\begin{bkfigure}[\linewidth]
\begin{minipage}{0.48\linewidth}\centering
\begin{tikzpicture}[scale=1.05]
  % Ratios measured off a 400 dpi crop of the book figure (book p.151):
  % Q at 0.369 and R at 0.620 of BC, apex almost exactly over Q so that AQ
  % reads as vertical.  A is nudged 0.04 left of Q so that angle AQC is
  % 91.3 deg, as the book draws it, rather than an unstated exact right angle.
  \coordinate (B) at (0,0);  \coordinate (C) at (3.45,0);
  \coordinate (A) at (1.2331,1.7664);
  \coordinate (Q) at ($(B)!0.369!(C)$);
  \coordinate (R) at ($(B)!0.620!(C)$);
  \draw[fig] (B) -- (C) -- (A) -- cycle;
  \draw[fig] (A) -- (Q);
  \draw[fig] (A) -- (R);
  \node[vlab,above] at (A) {$A$};
  \node[vlab,left]  at (B) {$B$};
  \node[vlab,right] at (C) {$C$};
  % the two angles at A carry arcs, as in the book
  \draw[tick] ($(A)+(-124.92:0.45)$) arc (-124.92:-88.70:0.45);
  \draw[tick] ($(A)+(-62.85:0.50)$)  arc (-62.85:-38.55:0.50);
  % each letter on the bisector of the angle it names, far enough out to clear
  % both sides of that angle
  \node[slab] at (0.4079,0.2127) {$x$};
  \node[slab] at (1.0307,1.0963) {$y$};
  \node[slab] at (1.7716,1.1087) {$s$};
  \node[slab] at (1.5527,0.2860) {$z$};
  \node[slab] at (1.7465,0.2398) {$t$};
  \node[slab] at (2.8459,0.2113) {$r$};
\end{tikzpicture}
\figcap{8--28}
\end{minipage}\hfill
\begin{minipage}{0.48\linewidth}\centering
\begin{tikzpicture}[scale=0.80]
  \coordinate (B) at (0,0);  \coordinate (C) at (3.05,0);
  % AB = BC: A on the circle centred B with radius BC, at 110 deg as drawn
  \coordinate (A) at ($(B)+(110:3.05)$);
  % D is where the 40-degree ray from B meets AC.  Solved exactly and written
  % out as literals so no extra tikz library is needed; angle CBD = 40.0000
  % deg is re-checked in tools/constraints/ch08.py.
  \coordinate (D) at (1.38738,1.16417);
  % BE = BD: E on BA at the same distance from B as D  (ratio 0.59381)
  \coordinate (E) at ($(B)!0.59381!(A)$);
  \draw[fig] (B) -- (C) -- (A) -- cycle;
  \draw[fig] (B) -- (D);
  \draw[fig] (E) -- (D);
  \node[vlab,above left] at (A) {$A$};
  \node[vlab,below] at (B) {$B$};
  \node[vlab,below right] at (C) {$C$};
  \node[vlab] at (1.6670,1.4040) {$D$};
  \node[vlab] at (-1.3700,1.4300) {$E$};
  \node[slab] at (1.2592,0.4583) {$40^\circ$};
\end{tikzpicture}
\figcap{8--29}
\end{minipage}
\end{bkfigure}}

% --------------------------------------------------------------- Fig 8-30
% AB congruent to A'B', with B between A and C.
\newcommand{\FIGVIIITHIRTY}{\begin{bkfigure}[0.66\linewidth]
\begin{tikzpicture}[scale=0.90]
  \draw[fig] (-0.20,0) -- (2.20,0);
  \coordinate (A) at (0,0); \coordinate (B) at (1.05,0);
  \coordinate (C) at (1.86,0);
  \foreach \p in {A,B,C}{\fill (\p) circle (1.0pt);}
  \draw[tick] ($(A)!0.5!(B)+(0,-0.13)$) -- ($(A)!0.5!(B)+(0,0.13)$);
  \node[vlab,above] at (A) {$A$};
  \node[vlab,above] at (B) {$B$};
  \node[vlab,above] at (C) {$C$};
  \begin{scope}[shift={(3.05,0)}]
    \draw[fig] (-0.20,0) -- (1.60,0);
    \coordinate (Ap) at (0,0); \coordinate (Bp) at (1.05,0);
    \fill (Ap) circle (1.0pt); \fill (Bp) circle (1.0pt);
    \draw[tick] ($(Ap)!0.5!(Bp)+(0,-0.13)$) -- ($(Ap)!0.5!(Bp)+(0,0.13)$);
    \node[vlab,above] at (Ap) {$A'$};
    \node[vlab,above] at (Bp) {$B'$};
  \end{scope}
\end{tikzpicture}
\figcap{8--30}
\end{bkfigure}}

% -------------------------------------------------------- Fig 8-31, 8-32
% 8-31: angles CAB and ABD right, CA = DB, M and N midpoints of CD and AB.
% 8-32: X and Y on one side of AB; rays AX and BY meet at P.
\newcommand{\FIGVIIITHIRTYONETWO}{\begin{bkfigure}[\linewidth]
\begin{minipage}{0.50\linewidth}\centering
\begin{tikzpicture}[scale=0.80]
  \coordinate (A) at (0,0);      \coordinate (B) at (3.05,0);
  \coordinate (C) at (0,1.62);   \coordinate (D) at (3.05,1.62);
  \coordinate (M) at ($(C)!0.5!(D)$);
  \coordinate (N) at ($(A)!0.5!(B)$);
  \draw[fig] (A) -- (B) -- (D) -- (C) -- cycle;
  \draw[aux] (M) -- (N);
  % right-angle boxes at A and B
  \draw[fig] ($(A)+(0,0.20)$) -- ($(A)+(0.20,0.20)$) -- ($(A)+(0.20,0)$);
  \draw[fig] ($(B)+(0,0.20)$) -- ($(B)+(-0.20,0.20)$) -- ($(B)+(-0.20,0)$);
  % single ticks on CM, MD ; double ticks on AN, NB
  \draw[tick] ($(C)!0.5!(M)+(0,-0.13)$) -- ($(C)!0.5!(M)+(0,0.13)$);
  \draw[tick] ($(M)!0.5!(D)+(0,-0.13)$) -- ($(M)!0.5!(D)+(0,0.13)$);
  \foreach \dx in {-0.06,0.06}{
    \draw[tick] ($(A)!0.5!(N)+(\dx,-0.13)$) -- ($(A)!0.5!(N)+(\dx,0.13)$);
    \draw[tick] ($(N)!0.5!(B)+(\dx,-0.13)$) -- ($(N)!0.5!(B)+(\dx,0.13)$);}
  \node[vlab,above left] at (C) {$C$};
  \node[vlab,above] at (M) {$M$};
  \node[vlab,above right] at (D) {$D$};
  \node[vlab,below left] at (A) {$A$};
  \node[vlab,below] at (N) {$N$};
  \node[vlab,below right] at (B) {$B$};
\end{tikzpicture}
\figcap{8--31}
\end{minipage}\hfill
\begin{minipage}{0.46\linewidth}\centering
\begin{tikzpicture}[scale=0.80]
  \coordinate (A) at (0,1.72);  \coordinate (B) at (0,0);
  \coordinate (P) at (3.05,1.02);
  \coordinate (X) at ($(A)!0.55!(P)$);
  \coordinate (Y) at ($(B)!0.55!(P)$);
  \draw[fig] (A) -- (B);
  \draw[fig] (A) -- (X);
  \draw[fig] (B) -- (Y);
  \draw[aux] (X) -- (P);
  \draw[aux] (Y) -- (P);
  \draw[fig] ($(P)+(0.17,0.20)$) -- (P) -- ($(P)+(0.17,-0.20)$);
  \fill (X) circle (1.0pt); \fill (Y) circle (1.0pt);
  \node[vlab,above] at (A) {$A$};
  \node[vlab,below] at (B) {$B$};
  \node[vlab,above] at (X) {$X$};
  \node[vlab,anchor=north] at ($(Y)+(0,-0.32)$) {$Y$};
  \node[vlab,above right] at (P) {$P$};
\end{tikzpicture}
\figcap{8--32}
\end{minipage}
\end{bkfigure}}
